\section{Conclusion}
\label{sec:conclusion}

This study shows that discrimination alone is not sufficient to establish the reliability of autism classification from resting-state functional connectomes. Using a fixed analytical framework, the effects of cohort composition, motion, regularisation, probability calibration, and model interpretation are examined alongside predictive performance. The ensemble achieves an AUROC of 0.776 and maintains similar discrimination when acquisition sites are withheld, but its performance varies with the cognitive composition of the autistic sample. Classifier scores separate diagnostic groups without showing a corrected association with symptom severity within autistic participants. Lower somatomotor connectivity also persists and strengthens after within-site matching on recorded mean FD, although this remains an exploratory group-level finding rather than evidence of a biomarker. At the modelling level, regularised linear classifiers perform better than the evaluated fixed-configuration tree ensembles, while the apparent advantage of windowed connectivity is not retained when covariance regularisation is matched. The native ensemble probabilities show a systematic calibration offset that is corrected using training-only recalibration without reducing \addhl{discrimination, and the offsets that remain between sites follow the autistic proportion of the site.} Interpretation is also more reliable at the level of functional systems: default-mode pairs contribute the largest overall share of classifier separation, while system-level attribution reproduces more consistently across cohort \addhl{halves, including halves that share no site,} than individual connection patterns. Taken together, these findings support evaluating connectome-based autism classifiers on criteria beyond discrimination, namely calibration, dependence on cohort composition, analytical design, and the reproducibility of model explanations. \addhl{Within ABIDE-I, the present ensemble meets several of these conditions: average calibration after training-only recalibration, a training-set threshold that holds at unseen sites, and system-level explanations that reproduce even between halves that share no site. Calibration at a new site requires its autistic proportion, and} independence from cohort composition and external validity are not established, so trustworthy connectome-based classification remains a goal rather than a result of this study. Independent external validation is required before these conclusions can be extended beyond the present retrospective research cohort.